\chapter*{英文术语与中文翻译规范}

\section*{英文术语保留规则}

本节旨在说明本书在术语使用与翻译上的基本原则，以便全书在行文、教学表达与术语对应上保持统一，也便于读者对照英文文献、论文和代码进行阅读。总体而言，我们采取“尊重学界通行用法、兼顾中文可读性与术语一致性”的原则：对于已经在深度学习与机器学习文献中固定使用、且在国际交流中具有高度识别性的英文术语，本书中文正文通常保留其英文拼写。具体而言，包括但不限于 Transformer、 token/tokens、 softmax 及常见模型缩写（如 BERT、 GPT、 CLIP、DDIM）。为便于读者理解，我们会在术语首次出现时酌情以括号补充简短中文说明（例如 token（数据元））。但在正文中，我们通常不以中文译名替代上述英文术语的规范写法，以保持与论文、开源代码和国际教材的一致性。


\section*{中英文翻译对照}

\begin{center}
\textbf{人工智能术语翻译对照}

\vspace{0.3em}

\begin{tabular}{ll}
\hline
\textbf{English} & \textbf{中文} \\
\hline
AI program & 人工智能纲领 \\
cybernetics program & 控制论纲领 \\
embodied AI & 具身智能 \\
imitation learning & 模仿学习 \\
ontogenetic & 个体成长智能 \\
open-ended & 开放式 \\
parsimony & 简约性 \\
representation & 表征 \\
representation learning & 表征学习 \\
scientific intelligence & 科学智能 \\
self-consistency & 自洽性 \\
societal intelligence & 社会智能 \\
\hline
\end{tabular}
\end{center}



\bigskip 

\begin{center}
\textbf{优化相关术语翻译对照}

\vspace{0.3em}

\begin{tabular}{ll}
\hline
\textbf{English} & \textbf{中文} \\
\hline
adjoint & 伴随 \\
alternating optimization & 交替优化 \\
automatic differentiation (AD) & 自动微分 \\
convex / strongly convex & 凸 / 强凸 \\
gradient descent & 梯度下降 \\
ill-conditioned & 条件不良的 \\
Kronecker factorization & 克罗内克分解 \\
learning rate & 学习率 \\
line search & 线搜索 \\
Nash equilibrium & 纳什均衡 \\
oracle & 预言机 \\
proximal gradient descent & 近端梯度下降 \\
scalable & 可扩展 \\
Stackelberg equilibrium & 斯塔克尔伯格均衡 \\
(sub)gradient & （次）梯度 \\
tractable & 易处理的（或可解的，视语境而定） \\
\hline
\end{tabular}
\end{center}

\bigskip
\begin{center}
\textbf{信息论与编码相关术语翻译对照}

\vspace{0.3em}

\begin{tabular}{ll}
\hline
\textbf{English} & \textbf{中文} \\
\hline
coding rate & 编码率 \\
coding rate reduction & 编码率降低 \\
compressive sensing & 压缩感知 \\
cross-entropy & 交叉熵 \\
differential entropy & 微分熵 \\
entropy & 熵 \\
Jensen--Shannon (JS) divergence & Jensen--Shannon 散度（JS 散度） \\
KL divergence & KL 散度 \\
Kolmogorov complexity & 柯尔莫哥洛夫复杂性 \\
lossy coding rate & 有损编码率 \\
mutual information & 互信息 \\
rate distortion & 失真编码率 \\
total variation (TV) distance & 总变差距离（TV 距离） \\
\hline
\end{tabular}
\end{center}

\bigskip
\begin{center}
\textbf{神经网络相关术语翻译对照}

\vspace{0.3em}

\begin{tabular}{ll}
\hline
\textbf{English} & \textbf{中文} \\
\hline
autoencoder (AE) & 自编码器 \\
Bayes optimal denoiser & 贝叶斯最优去噪器 \\
classifier guidance & 分类器引导 \\
classifier-free guidance (CFG) & 无分类器引导 \\
complete/overcomplete & 完备/过完备 \\
consistency model & 一致性模型 \\
consistent representation & 一致表征 \\
cross attention & 交叉注意力（写法全书统一） \\
denoising & 去噪 \\
diffusion process / diffusion models & 扩散过程 / 扩散模型 \\
embedding & 嵌入 \\
empirical distribution & 经验分布 \\
few-step model & 少步模型 \\
flow matching & 流匹配 \\
generalization & 泛化 \\
ground truth & 真实值（或真值） \\
latent (space) & 潜在（空间） \\
memorization & 记忆 \\
noise predictor & 噪声预测器 \\
representation autoencoder (RAE) & 表征自编码器 \\
sampling & 采样 \\
text prompt & 文本提示 \\
transport map & 传输映射 \\
velocity predictor & 速度预测器 \\
\hline
\end{tabular}
\end{center}



